% 图契约：本题的源、边界和推导逐项对应；箭头符号按图注解释。
\documentclass[tikz,border=7pt]{standalone}
\usepackage{fontspec,xeCJK,amsmath,amssymb}
\IfFontExistsTF{Times New Roman}{\setmainfont{Times New Roman}}{\setmainfont{TeX Gyre Termes}}
\IfFontExistsTF{Heiti SC}{\setCJKmainfont[AutoFakeBold=1.5]{Heiti SC}}{\setCJKmainfont[AutoFakeBold=1.5]{FandolHei-Regular}}
\usetikzlibrary{arrows.meta,bending,calc,patterns,positioning}
\definecolor{ink}{HTML}{243745}
\definecolor{blue}{HTML}{275DAB}
\definecolor{green}{HTML}{177C63}
\definecolor{red}{HTML}{BD4935}
\definecolor{orange}{HTML}{AA6C16}
\definecolor{purple}{HTML}{715393}
\tikzset{every node/.style={font=\fontsize{9.3}{12.5}\selectfont,text=ink},
axis/.style={-{Stealth[length=4pt]},line width=.65pt,draw=ink!65},
vec/.style={-{Stealth[length=5pt]},line width=1pt,draw=blue},
short/.style={-{Stealth[length=3pt]},line width=.8pt},
guide/.style={dashed,draw=ink!40,line width=.5pt},
card/.style={rounded corners=2pt,draw=blue!65,fill=blue!4,align=center,text width=6.5cm,inner xsep=9pt,inner ysep=7pt},
title/.style={font=\fontsize{10.5}{14}\selectfont\bfseries,text width=15cm,align=center},
label/.style={font=\fontsize{8.4}{11}\selectfont,align=center},
note/.style={label,text width=6.2cm}}
\begin{document}
\begin{tikzpicture}[x=1cm,y=1cm]
\path[use as bounding box] (-8,-3.5) rectangle(8,3.5);
\node[title] at (0,3.15) {M2-176c\quad 窄隙中的局部厚度决定对数主项};
\begin{scope}[shift={(-4,0)}]
\draw[axis](-2.4,0)--(2.4,0)node[right]{$\rho$};
\draw[blue,thick]plot[domain=-2:2] (\x,{.22+.19*\x*\x});\draw[green,thick]plot[domain=-2:2](\x,{-.22-.19*\x*\x});
\draw[<->,orange](0,-.22)--(0,.22);\node[above right]at(.1,.25){$\Delta$};
\draw[<->,orange](1.5,-.6475)--(1.5,.6475);\node[fill=white,inner sep=2pt]at(1.5,1.3){$h(\rho)$};
\node[note]at(0,-1.8){同心薄环面积 $dA=2\pi\rho\,d\rho$；\\局部模型只确定对数的系数。};
\end{scope}\node[card,draw=blue!65,fill=blue!4]at(3.9,1.8){$h(\rho)\simeq\Delta+\rho^2/R$};
\node[card,draw=green!65,fill=green!4]at(3.9,0){$dC\simeq\varepsilon_0\,2\pi\rho\,d\rho/h(\rho)$};
\node[card,draw=orange!65,fill=orange!4]at(3.9,-1.8){$C_{\text{pair}}\sim\pi\varepsilon_0R[\ln(4R/\Delta)+2\gamma]$};
\end{tikzpicture}\end{document}